% cv_body.tex — CV content; included by cv.tex (layout chosen there)

% ---------- Header ----------
\begin{center}
  \cvname{Joon-Hyuk Ko}\\[0.5em]
  Center for AI and Natural Sciences, Korea Institute for Advanced Study (KIAS)\\
  85 Hoegi-ro, Dongdaemun-gu, Seoul 02455, Republic of Korea\\[0.35em]
  \href{mailto:jhko725@kias.re.kr}{jhko725@kias.re.kr} \enspace$\cdot$\enspace
  +82\,10-8919-1384 \enspace$\cdot$\enspace
  \href{https://scholar.google.com/citations?user=t7lTWNQAAAAJ&hl=en}{Google Scholar} \enspace$\cdot$\enspace
  \href{https://github.com/Jhko725}{GitHub}
\end{center}

% ---------- Research Interests ----------
\section{Research Interests}
Scientific machine learning for dynamical systems: neural surrogates (neural ODEs, neural operators) whose training and architecture are grounded in dynamical systems theory and numerical analysis.
\begin{interests}
  \item \textbf{Chaotic systems:} training with synchronization and local phase-space geometry to recover correct long-term statistics
  \item \textbf{Parametrized systems:} generalization across parameters and bifurcations via numerical continuation
  \item \textbf{Conservation laws:} in-context neural operators built on conservative finite-volume structure
  \item \textbf{Oscillators:} phase--amplitude reduction through normal forms and learned conjugacies
  \item \textbf{Applications:} inverse problems in experimental physics, especially high-speed atomic force microscopy and microscopy of soft and living matter
\end{interests}

% ---------- Positions ----------
\section{Positions}
\entry{AI Research Fellow}{2025 -- present}{Center for AI and Natural Sciences, Korea Institute for Advanced Study, Seoul}

% ---------- Education ----------
\section{Education}
\entry{Ph.D.\ in Physics}{2019 -- 2025}{Seoul National University. Advisor: Prof.\ Wonho Jhe}
\begin{notes}
  \item Dissertation: \textit{Inverse problems in atomic force microscopy: Universal force reconstruction and viscoelasticity extraction}
\end{notes}
\entry{B.A.\ in Physics, \textit{summa cum laude}}{2016 -- 2019}{Seoul National University. Early graduation in three years}
\entry{Korea Science Academy of KAIST}{2013 -- 2016}{\relax}

% ---------- Publications ----------
\section{Publications}
% publications.tex — shared between the CV and the publication list
% \pub{title}{authors}{venue};  \cofirst = co-first author, \cocorr = co-corresponding author

{\small\color{softgray} \cofirst~Co-first author (equal contribution);\enspace \cocorr~Co-corresponding author.}

\subsection{Journal and conference publications}
\begin{pubs}
  \pub{Desmin filaments follow contractile actin in muscle sarcomere interior.}
      {S.~Joseph\cofirst, \me\cofirst, N.~Dahan, T.~Dadosh, H.~K.~Pak, J.~A.~Kirk, W.~Jhe\cocorr, and S.~Shemer\cocorr}
      {\textit{Nature Communications} (2026), in press.}
  \pub{Quartz resonator-based atomic force microscopy as a stable semiconductor metrology solution.}
      {H.~Koh, J.~Seo, J.~Yun, H.~Kang, H.~Lee, S.~Do, H.~Choe, \me, and W.~Jhe}
      {\textit{Metrology, Inspection, and Process Control XL}, Proc.\ SPIE \textbf{13981} (2026).}
  \pub{Homotopy-based training of NeuralODEs for accurate dynamics discovery.}
      {\me\cofirst, H.~Koh\cofirst, N.~Park, and W.~Jhe}
      {\textit{Advances in Neural Information Processing Systems (NeurIPS)} \textbf{36} (2023).}
  \pub{Bayesian neural network with pretrained protein embedding enhances prediction of drug--protein interaction.}
      {Q.~H.~Kim, \me, S.~Kim, N.~Park, and W.~Jhe}
      {\textit{Bioinformatics} \textbf{37}(20) (2021).}
  \pub{A semiautomated measurement of muscle fiber size using the Imaris software.}
      {J.~E.~Gilda\cofirst, \me\cofirst, A.~Y.~Elfassy, N.~Tropp, A.~Parnis, B.~Ayalon, W.~Jhe, and S.~Cohen}
      {\textit{American Journal of Physiology -- Cell Physiology} \textbf{321}(3) (2021).}
  \pub{Universal theory of dynamic force microscopy for exact and robust force reconstruction using multiharmonic signal analysis.}
      {S.~Kim, \me, and W.~Jhe}
      {\textit{Physical Review Letters} \textbf{126}(7) (2021).}
  \pub{GCIceNet: a graph convolutional network for accurate classification of water phases.}
      {Q.~H.~Kim, \me, S.~Kim, and W.~Jhe}
      {\textit{Physical Chemistry Chemical Physics} \textbf{22}(45) (2020).}
\end{pubs}

\subsection{Preprints}
\begin{pubs}
  \pub{Watch your neighbors: Training statistically accurate chaotic systems with local phase space information.}
      {\me, A.~Giraldo, and D.-S.~Lee}
      {\href{https://arxiv.org/abs/2605.14405}{arXiv:2605.14405} (2026). Under review at ICLR 2027.}
  \pub{A robust foundation model for conservation laws: Injecting context into flux neural operators via recurrent vision transformers.}
      {T.~Kim\cofirst\ and \me\cofirst}
      {\href{https://arxiv.org/abs/2605.05488}{arXiv:2605.05488} (2026). Under review at ICLR 2027.}
  \pub{A comprehensive semi-automated fabrication system for quartz tuning fork AFM probes with real-time resonance frequency monitoring and Q-factor control.}
      {H.~Koh, \me, and W.~Jhe}
      {\href{https://arxiv.org/abs/2601.16221}{arXiv:2601.16221} (2026).}
\end{pubs}


% ---------- Talks ----------
\section{Talks}
\subsection{Invited talks and lectures}
\begin{talks}
  \talk{Aerospace engineering $\times$ atomic force microscopy: Opportunities for scientific machine learning and control}
       {Aerospace Autonomy and Intelligence Laboratory, Dept.\ of Aerospace Engineering, Seoul National University}{Jul.\ 2026}
  \talk{Applying AI to experimental physics research---Case study: Data-driven force reconstruction in atomic force microscopy}
       {Advanced Topics in Applied Physics (graduate course), Dept.\ of Physics, Seoul National University}{May 2026}
  \talk{Inverse problems in atomic force microscopy}
       {Surface Physics (graduate course), Dept.\ of Physics, Jeonbuk National University}{May 2026}
  \talk{Data-driven dynamics discovery with neural ordinary differential equations}
       {Innovation Center for Industrial Mathematics, National Institute for Mathematical Sciences}{Mar.\ 2026}
  \talk{Data-driven dynamics discovery with neural ordinary differential equations}
       {National Institute for Mathematical Sciences}{Feb.\ 2026}
  \talk{Effectively training neural ordinary differential equations for data-driven dynamics discovery}
       {Dept.\ of Mathematical Sciences Seminar, Seoul National University}{Jan.\ 2026}
\end{talks}

\subsection{Conference and workshop talks}
\begin{talks}
  \talk{Learning phase--amplitude coordinates of oscillators via neural network conjugacy to a normal form}
       {Korean Society for Industrial and Applied Mathematics (KSIAM) 2026 Fall Meeting}{Nov.\ 2026}
  \talk{Training statistically accurate dynamical models with local phase space information}
       {KSIAM 2026 Spring Meeting}{May 2026}
  \talk{Effective training of neural ordinary differential equations for data-driven discovery of chaotic dynamics}
       {PDE and Applied Mathematics Winter Workshop 2026}{Jan.\ 2026}
  \talk{Learning statistically accurate neural surrogate dynamics for chaotic systems through phase space neighborhood information}
       {KIAS Center for AI and Natural Sciences Winter Workshop}{Jan.\ 2026}
  \talk{Neighborhood-aware training of NeuralODEs to accurately predict dynamical invariants of chaotic systems}
       {KIAS Center for AI and Natural Sciences Workshop}{May 2025}
  \talk{On the difficulty of training autoregressive models on long time series data}
       {KIAS Computational Sciences Workshop}{May 2025}
  \talk{Scientific machine learning with neural ordinary differential equations}
       {KIAS Center for AI and Natural Sciences AI Seminar}{Apr.\ 2025}
\end{talks}

\subsection{Poster}
\begin{talks}
  \talk{Neural network extraction of viscoelasticity from AFM indentation measurements\newline
        {\upshape S.~Hong, \me, H.~Koh, and W.~Jhe}}
       {7th International Soft Matter Conference, Osaka, Japan}{Sep.\ 2023}
\end{talks}

% ---------- Awards ----------
\section{Awards and Honors}
\begin{dated}
  \datedItem{\textbf{Talent Award of Korea}, Korean Ministry of Education (100 awardees nationwide per year)}{2015}
  \datedItem{\textbf{4th place, International Young Physicists' Tournament}, Korean national team}{2014}
\end{dated}

% ---------- Service & Outreach ----------
\section{Service and Outreach}
\subsection{Reviewing}
Neural Information Processing Systems (NeurIPS); International Conference on Learning Representations (ICLR); IEEE Transactions on Pattern Analysis and Machine Intelligence.

\subsection{Public lectures and media}
\begin{dated}
  \datedItem{Interview with KBS News on students' use of AI}{Feb.\ 2026}
  \datedItem{\textit{Studying smart with AI}. Public lecture, Dongdaemun-gu $\times$ KIAS AI, Math, and Science Academy}{Feb.\ 2026}
  \datedItem{Introduced the Center for AI and Natural Sciences in the KIAS interview series}{Dec.\ 2025}
  \datedItem{\textit{How does AI work / Using AI in everyday life}. Public lecture, KIAS Summer Science Academy}{Aug.\ 2025}
\end{dated}

% ---------- Research Experience ----------
\section{Additional Research Experience}
\entry{Visiting Fellow, Harvard University}{2020 -- 2022}{Experimental Soft Condensed Matter Group (Prof.\ David A.\ Weitz), SEAS}
\begin{notes}
  \item Built a quartz tuning fork AFM for a joint Jhe--Weitz project; machine learning-based viscoelasticity extraction for soft materials. Began the collaboration with Prof.\ Shenhav Shemer (Technion) behind publications [1] and [5].
\end{notes}
\entry{Undergraduate Research Intern, Seoul National University}{2017 -- 2019}{Jhe Lab, Dept.\ of Physics. Built a quartz tuning fork AFM and its LabVIEW control software.}
